% Generated from validated Article Draft Result. Do not hand-edit; revise the structured draft instead.

OpenAI o1のpreview、ChatGPT上の製品展開、API availabilityは一つのreleaseとして潰さず、異なるlifecycle eventとして追う。同時期のQwQやDeepSeekのreasoning-oriented previewは、各社が推論時の計算配分を新しい能力軸として試したことを示すが、名称が似ていても同一のmechanismや同一のavailabilityを意味しない。\autocite{src-14a855188533,src-a452b2dc7f5d,src-ea823ed75f72}


研究側では、test-time computeを増やすことで探索・検証・解答選択を改善しようとする試みと、それを測る難しい評価系が並行して現れた。ここで重要なのは、test-time scalingという一般的な設計空間と、特定vendorが『reasoning model』として提供する製品を同義にしないことである。評価benchmarkやauthor-reported resultは、設定と主張主体を保ったまま読む。\autocite{src-1f00c386f820,src-259b57a1bad9,src-7c2651c80f45,src-9cf289b414fd}


\begin{claimboundary}
2024年末時点で観測できたreasoningの潮流を、2025年に公開されたDeepSeek-R1の結果から逆算して説明しない。2024-H2では当時のpreview、release、研究結果だけを根拠に、推論時計算が独立した設計軸として立ち上がったところまでを結論とする。\autocite{src-ea823ed75f72}
\end{claimboundary}
